\documentclass[../report_main.tex]{subfiles}
\begin{document}

\section{Implemented Functionalities}
The proposed database architecture successfully realizes a comprehensive suite of functionalities required for an enterprise-level online examination platform. Foremost, it enforces a mathematically sound, role-based access control mechanism that securely distinguishes between administrative, instructional, and student entities. Educators are provided with the structural capability to manage dynamic classroom rosters and construct highly polymorphic question banks. The system facilitates the execution of standardized assessments utilizing flexible, JSON-based scoring configurations. Furthermore, the architecture effectively captures granular transaction logs of student responses, leveraging database triggers and stored procedures to automate the aggregation of final performance metrics with minimal application-layer intervention.

\section{Architectural Challenges and Solutions}
Throughout the conceptualization and physical design phases, several prominent structural challenges were encountered. A primary difficulty involved engineering a repository capable of storing both atomic questions and complex, hierarchical item sets (e.g., reading comprehension passages) without violating normalization axioms. This complication was resolved by introducing a self-referencing foreign key constraint within the primary \texttt{questions} table, which accommodates infinite hierarchical depth. Additionally, addressing the diverse attributes inherent to distinct question typologies presented a risk of schema sparsity. The implementation of a polymorphic design pattern-retaining generalized metadata in a central hub while delegating format-specific constraints to auxiliary tables-effectively eliminated null-value proliferation and ensured optimal storage density.

\section{Evaluation of the Final Architecture}
The final database design offers several key advantages. Its modular structure makes it highly scalable, allowing new question formats or media to be added easily without disrupting the existing system. Additionally, strict foreign key rules prevent orphan records, ensuring that student test histories remain consistent and accurate. However, there are some trade-offs. Because of the highly modular and polymorphic design used to store various question types, retrieving a complete examination often requires complex joins across multiple tables. While B-Tree indexes currently help speed up these queries, future versions of the system could use materialized views to pre-calculate complex data for faster reporting.

\section{Collaborative Contributions}
The successful delivery of this database engineering project was achieved through structured, collaborative efforts across multiple developmental domains. The workload was systematically partitioned to ensure academic rigor at every phase. The project encompassed the drafting of conceptual Entity-Relationship models, the translation of logical schemas into optimized Data Definition Language (DDL) scripts, and the establishment of complex procedural logic. The deployment of advanced SQL formulations, performance indexing, and rigorous documentation collectively ensured that the final system satisfied all technical requirements and academic standards expected of a sophisticated relational database implementation.

\end{document}